\chapter*{摘要}
\addcontentsline{toc}{chapter}{摘要}
本项目围绕视觉语言模型参与机器人操作的实际作用，逐步完成桌面原型、公开策略复现、高层任务处理、视觉完成检查、同场景移动操作、可供性定位以及端侧部署。早期先建立SO101机械臂的遥操作与观察工具，再通过动作分块策略和真实示教检查数据与动作接口；随后收敛到LIBERO公开操作基准，使用\pifive{}模型建立稳定的视觉到动作执行基线。

在基线上，项目检验Qwen对中文指令分解、阶段检查和失败恢复的作用。对照结果表明，语言翻译、视觉判断和动作执行是不同能力，高层给出正确目标也不保证低层策略按目标行动。后续将定位问题收敛到抽屉把手等可操作位置，训练低秩适配器，并接入同一物理世界中的导航、规划和实际模拟关节控制。端侧则在独立Jetson环境中比较小模型、量化精度、设备输出和推理时延。

目前，\pifive{}公开任务基线为117/120；冻结中文分层版本为28/30，仅翻译对照为29/30。最新可供性驱动预接触测试中，LoRA版本为4/4，基座为0/4；额外前置姿态后的完整任务交接为2/4，说明接口贯通仍需技能阶段兼容性验证。Jetson感知与PC仿真执行的三个预接触回合通过，但15瓦模式下的时延尚未达到单次3秒目标。这些数字来自不同任务和条件，不能相加成为一个总体成功率。

报告按六个技术篇章、二十四个阶段组织，保留从早期理想条件实验到公开基准、再到同场景物理执行的演进过程。实验设置、数据划分、主要对照、失败情形和代码位置在相应章节展开；复现命令与素材索引集中列于附录。历史方案与修正后的方案分别标识，尤其将动作适配器的权重合并问题与策略本身的能力分开讨论。

研究结果支持三个工程判断：首先，需要用仅翻译、固定切换和原生策略等简单对照判断复杂模块是否有实际贡献；其次，二维定位、三维反投影、规划到位与完整任务成功应分层验收；最后，端侧能加载模型不等于满足实时要求。当前可复现的主要贡献是系统集成、受控对照与接口诊断，还不能据此声称具备跨场景通用操作能力。

\noindent\textbf{关键词：}视觉语言动作模型；机器人操作；可供性定位；同场景仿真；低秩适配；端侧部署。
